\section{Limits, boundary traces, and absolute-height identities}
\label{sec:limits}

We continue to use the integer height $H_\delta=(h_\delta-1)/2$,
and write $H_\delta(f)=\int H_\delta(x)f(x)\,dx$. All tests in this
section are supported in an eventual interior compact of $D$. Thus an
arbitrary extension from $D_\delta$ to $D$, for example extension by zero,
has no effect on the limits under consideration. A bracket means one of
the two filled-window comparison laws of
\cref{prop:bracket-identity}. Its boundary trace, although random, takes
values in $[-2,1]$.

\begin{proposition}[Tightness and convergence of moments]
\label{prop:tightness}
The true fields and all the bracket fields are tight in
$H^{-t}_{\mathrm{loc}}(D)$ for every $t>1$. For $K\Subset D$, $p<\infty$,
and tests supported in $K$,
\begin{equation}
 \sup_\delta\norm{H_\delta(f)}_{L^p}
       \le C_{K,p}\norm{f}_{\infty}.
 \label{eq:local-test-moments}
\end{equation}
The constants are uniform over these bracket laws. Along convergence in
law, every finite joint test moment converges. These assertions hold
also for any fixed finite collection of fields in a coupling.
\end{proposition}

\begin{proof}
Cover $K$ by finitely many balls whose radii are a sufficiently small
fixed fraction of their distance from $\partial D$, and use a partition
of unity. The smearing estimate of \cref{prop:smearing-moments}, applied
to the positive and negative parts on each ball, gives
\eqref{eq:local-test-moments}. For brackets, condition first on the
boundary trace. Its uniform bound gives the same constants before and
after averaging over that trace.

For completeness, place a cutoff $\chi\in C_c^\infty(D)$ in a fixed
square, regarded as a flat torus, and let $e_n$, $n\in\Z^2$, be its
normalized Fourier basis. Since $\norm{\chi e_n}_\infty$ is bounded
independently of $n$, for $s>1$ we have
\begin{equation}
 \sup_\delta\EE\norm{\chi H_\delta}_{H^{-s}}^2
 =\sup_\delta\sum_{n\in\Z^2}(1+|n|^2)^{-s}
                  \EE|H_\delta(\chi\overline{e_n})|^2<\infty.
 \label{eq:sobolev-bound}
\end{equation}
For any $p\ge2$, Minkowski's inequality applied to the sum defining
the squared norm gives its $L^{p/2}$ bound as well, using
\eqref{eq:local-test-moments} at order $p$. If $1<s<t$, the inclusion
$H^{-s}\hookrightarrow H^{-t}$ on this torus is compact: its Fourier
multiplier tends to zero, and truncating the Fourier series gives
finite-rank approximations. Markov's inequality and
\eqref{eq:sobolev-bound} therefore give tightness in $H^{-t}$.
Use a nested countable family of cutoffs equal to one on an exhaustion
of $D$. Choosing the compact-containment probabilities for its members
with a summable error gives tightness in the corresponding local
Sobolev space, and hence in distributions.

If $f_1,\ldots,f_j$ are fixed tests, H\"older's inequality and
\eqref{eq:local-test-moments} at order $j(1+\eta)$ bound
$\prod_i H_\delta(f_i)$ in $L^{1+\eta}$, for any fixed $\eta>0$.
It is uniformly integrable. Convergence in law of the test vector
therefore implies convergence of its product expectation. The same
argument applies to mixed products from a finite coupling. The
moment inequalities themselves pass to every limit by truncation.
\end{proof}

Fix once and for all a nonnegative radial function
$\rho\in C_c^\infty(B(0,1))$ with integral one, and set
\[
 \rho_{x,r}(z)=r^{-2}\rho((z-x)/r),\qquad
 d(x)=\dist(x,\partial D).
\]
Choose a sufficiently small fixed $\kappa>0$. The smearing estimate,
first with mesh tending to zero at fixed $x$, also implies
\begin{equation}
 \sup_{x\in D}\norm{H(\rho_{x,\kappa d(x)})}_{L^p}\le C_p
 \label{eq:depth-moment-bound}
\end{equation}
for any subsequential true or bracket limit. Here and below a harmless
change of the constant covers centers in a fixed compact part of $D$.
In particular this bound is valid when the center approaches a marked
point; it requires only bounded boundary data, not a constant arc.

\subsection{The ordered limits have the correct boundary mean}

We specify the limits of brackets before using their identities.
Choose increasing finite grids $\mathcal G_n$ whose union is dense
among the transverse regular horizontal levels of $\partial D$ and
which miss the two marks. Each such level has finitely many boundary
intersections. For fixed $n$, choose disjoint windows at these
intersections, avoiding the marks, and denote their maximum size by
$\eta$. At each mesh, height comparison gives a coupling
\begin{equation}
                 H_{\delta,n,\eta}^-
                    \le H_\delta\le H_{\delta,n,\eta}^+.
 \label{eq:ordered-brackets}
\end{equation}
For example, couple each ordered pair and disintegrate over the true
field to join the couplings. Affine interpolation preserves the order.

Start with an arbitrary convergent subsequence of the true laws, with
limit denoted by $H$. For each fixed $n,\eta$ take further joint
subsequences in \eqref{eq:ordered-brackets}; then let $\eta$ decrease
to zero at fixed $n$; finally let $n$ increase. Tightness permits all
these extractions. At every stage the true marginal is the originally
chosen law of $H$. The order persists when tested against nonnegative
smooth functions, because that is a closed condition under
distributional convergence. The use of a countable determining family
of nonnegative tests makes the order simultaneous almost surely.
Write the final triple as $(H^-,H,H^+)$.

Identities from \cref{prop:bracket-identity} pass through these limits.
For bounded continuous functions of local increment tests this follows
from convergence in law and the moment bound for the Laplace
insertion. Every admissible finite selection of grid levels belongs
to all sufficiently late grids. Taking the other insertion to be one
shows that the distributional means $m^-$ and $m^+$ of the final
brackets are harmonic locally throughout $D$. Indeed any interior
point admits two regular levels on opposite sides of a sufficiently
small testing disk. Distributional harmonic functions are smooth,
so these means are ordinary harmonic functions. Comparison with
constant-boundary laws, whose means equal the boundary constant by
height reflection, gives
\begin{equation}
                     -2\le m^-\le m^+\le1.
 \label{eq:bounded-bracket-means}
\end{equation}

\begin{proposition}[Boundary trace and mixed boundary moments]
\label{prop:boundary-trace}
Both limiting bracket means are the harmonic extension $m$ of $g$.
Consequently $H^-=H=H^+$ almost surely in the ordered coupling, and
$\EE H=m$.

Let $p$ be on an open boundary arc on which $g=k$, where $k$ is
$0$ or $-1$, and let $x\to p$ from within $D$. For the true limiting
field put $Y_x=H(\rho_{x,\kappa d(x)})$. If $F$ is a product of a fixed
finite number of height averages supported away from a neighborhood
of $p$, then
\begin{equation}
                      \EE[(Y_x-k)F]\longrightarrow0.
 \label{eq:mixed-boundary-trace}
\end{equation}
The same assertion holds when those averages approach finitely many
other, distinct, nonmarked boundary points, with their radii a fixed
small fraction of their depths. In that assertion the points and their
finite number are fixed before the depths tend to zero.
\end{proposition}

\begin{proof}
We first identify the bracket means. Fix a point $p$ on an open arc
of value $k$, a center $x$ near $p$, and its testing ball of radius
$\kappa d(x)$. Choose finitely many final arch scales surrounding
this ball in a small clean neighborhood of the nearest boundary
projection of $x$. Work initially at finite mesh with a fixed grid
and windows sufficiently small for the construction below, and write
$Y_{\delta,x}$ for the chosen bracket's unit average in the ball.
Two searches will isolate the ball from the exceptional window values
and then from the remote boundary data.

For each window choose a clean censor neighborhood on its constant
arc, of value $k_i$. Choose these neighborhoods disjoint, away from
the marks, small relative to the final arch scales, and small enough
to miss the testing ball. Conditional
on the original bracket trace, all its values are at least $k_i-2$,
and those in this censor neighborhood outside its window are at least
$k_i$. Enclose the entire discrete window and its adjacent sites in
the inner comparison box, with a positive margin before sending the
mesh to zero. Theorem~\ref{thm:boundary-arches} supplies an arch at
level $k_i$ surrounding this box. For a fixed grid there are only
finitely many windows; shrinking them relative to their censor
neighborhoods makes the total failure probability as small as desired.

On success choose the innermost arch by exploring from the window
side, as in \cref{lem:extremal-stopping}. Discard the patch between
the window and this arch, and retain the side away from the window.
This retained side contains the testing ball and has not been
examined by the search. Conditional on the revealed patches and
separator traces, its free heights have their Gibbs law. No original
exceptional vertex is adjacent to a retained free site. The new usable
separator sites carry their actual heights, which are at least $k_i$;
the global lower bounds are also preserved. Thus the first searches
replace the windows by small boundary pieces with the required
one-sided data. They do not condition on the preliminary revelations
used only to prove the arch probability estimate.

Apply \cref{cor:boundary-robustness} at the chosen final scales in
the retained domain. The censor patches are small enough for its
replacement hypothesis. In the chosen clean neighborhood every usable
original or replacement boundary site has height at least $k$.
With high probability there is an arch at level $k$ separating the
testing ball from the remote boundary. Choose the outermost such
arch by searching from the remote side. This time the unexamined
region is the side containing the testing ball, between the final
arch and the clean boundary arc. Its usable boundary consists of
sites on the arch and any original or censor pieces it contains,
all at height at least $k$. Height comparison with constant data
therefore gives conditional mean at least $k$ for $Y_{\delta,x}$.

Let $E_\delta$ denote success of both stages, and expose only the
discarded patches, the remote side of the final arch, and their
usable traces. On $E_\delta$ this information determines the event
and leaves the testing region Gibbs. For any $q>1$ the preceding
conditional comparison gives
\begin{equation}
 \EE(Y_{\delta,x}-k)\ge
       -\norm{Y_{\delta,x}-k}_{L^q}
                           \PP(E_\delta^c)^{1-1/q}.
 \label{eq:boundary-mean-error}
\end{equation}
The norm is bounded by the unconditional comparable-depth smearing
estimate, whose limit form is \eqref{eq:depth-moment-bound}. No
moment estimate after conditioning on censor success is needed.

For the upper bound start afresh from the original bracket trace.
Its global values are at most $k_i+2$, and its clean values in the
$i$th censor neighborhood are at most $k_i$. Use decreasing censor
arches and then a decreasing final arch in the same two stopping
directions. The testing region now has usable boundary at most $k$,
giving the matching upper estimate, with the same H\"older bound on
the failure contribution. This separate
construction is necessary because the lower censor arches carry no
upper bound on their revealed heights.

We make explicit the uniformity that permits the grid limit. Given an
error tolerance, first choose finitely many final arch blocks in a
sufficiently straight part of the limiting smooth arc. They are
available for every sufficiently shallow testing center. For each
fixed grid, choose the outer censor sizes so small that the robust
comparison applies in these blocks; the synthetic interval endpoints
in that comparison avoid all the tiny replacements. By
\cref{cor:boundary-robustness}, the final block estimates do not incur
a factor for the number of replacements. Next decrease the window
sizes so that all censor searches succeed with the desired total
error, and before that take the mesh and boundary approximation
errors to zero. The resulting bound on the mean depends on the final
blocks and the error tolerance, and not on the grid. It thus survives
window limits followed by grid refinement. This is exactly the order
of limits used to define $(H^-,H,H^+)$. Nearest projections and the
uniform smooth-arc estimates allow tangential as well as normal
approaches to $p$.

Since each bracket mean is harmonic, radial averaging recovers its
value at $x$. The two bounds just proved show that its boundary trace
is $k$ on the open arc. Together with
\eqref{eq:bounded-bracket-means}, bounded harmonic uniqueness gives
$m^-=m^+=m$. The two marked points do not change this uniqueness:
for a bounded harmonic function with zero trace elsewhere, compare
its absolute value with its supremum times the harmonic measure of
arbitrarily small boundary neighborhoods of the marks. At any fixed
interior point these harmonic measures tend to zero, since a smooth
domain has a Poisson kernel and single points have zero arclength.

The order now identifies the fields, rather than merely their means.
For every nonnegative test $f$,
$H^+(f)-H^-(f)\ge0$ and its expectation is zero. Thus it vanishes
almost surely. Apply this to a countable dense family of nonnegative
tests and use continuity of distributions. The ordered fields agree
as distributions almost surely, proving the first assertion and
transferring the bracket identities to $H$.

For the mixed assertion no censoring is needed in the true law.
Choose deterministic outer arch neighborhoods that do not meet the
supports of the other observations. First suppose $F\ge0$ is
measurable outside those neighborhoods. On the lower good event
$E_-$, choose the outermost lower arch from the exterior. The event,
the selected trace, and $F$ are measurable in the revealed exterior;
the Gibbs property gives
\[
                  \EE[(Y_x-k)F;E_-]\ge0.
\]
The upper search separately gives an event $E_+$ on which the reverse
inequality holds. These two searches need not share a sigma-field.
For $q>1$, their failure contributions satisfy
\begin{equation}
 -\norm{(Y_x-k)F}_{L^q}\PP(E_-^c)^{1-1/q}
 \le \EE[(Y_x-k)F]
 \le \norm{(Y_x-k)F}_{L^q}\PP(E_+^c)^{1-1/q}.
 \label{eq:weighted-arch-bounds}
\end{equation}
All norms here are bounded by H\"older and the smearing estimates.
For signed $F$, apply the bounds to its positive and negative parts;
they are exterior measurable and have the same moment control.
One uses the actual exterior stopping search at this step, without
retaining any of the preliminary revelations used to prove the arch
probability estimate. The result follows first at finite mesh, then
for limit moments by \cref{prop:tightness} and truncation.

When other observations approach different boundary points, fix that
finite set first and take the outer neighborhoods smaller than their
mutual distances. The observations remain exterior measurable as
their depths decrease. Estimate \eqref{eq:depth-moment-bound} bounds
their moments uniformly, so the same proof gives
\eqref{eq:mixed-boundary-trace} in this joint limiting regime.
\end{proof}

\begin{lemma}[Deterministic boundary reference averages]
\label{lem:boundary-references}
Fix a compact subarc of one open constant boundary arc, with value
$k$. There exist finite averages of comparable-depth unit smearings
near that subarc which converge to $k$ in every finite $L^p$ of the
true limiting field. More precisely, choose $N$ distinct points
$p_1,\ldots,p_N$ in the subarc, let $x_i(t)$ approach $p_i$ normally
at depth $t$, and put
\[
 A_{N,t}=\frac1N\sum_{i=1}^N
       H(\rho_{x_i(t),\kappa d(x_i(t))}).
\]
For all sufficiently small $t$ the balls are disjoint, and
\begin{equation}
 \limsup_{t\downarrow0}\EE|A_{N,t}-k|^2\le\frac C N,
 \qquad
 \sup_{N,t}\norm{A_{N,t}-k}_{L^q}\le C_q.
 \label{eq:boundary-average-variance}
\end{equation}
The second supremum is over admissible small depths. Consequently
$\lim_{N\to\infty}\limsup_{t\downarrow0}
\norm{A_{N,t}-k}_{L^p}=0$ for every finite $p$.
\end{lemma}

\begin{proof}
Write $Y_i(t)$ for the $i$th summand. Uniform comparable-depth
moments give $\EE|Y_i(t)-k|^2\le C$. For $i\ne l$, apply
\eqref{eq:mixed-boundary-trace} with exterior factor $Y_l(t)-k$:
\[
       \EE[(Y_i(t)-k)(Y_l(t)-k)]\longrightarrow0.
\]
Expanding the square, taking the depth limit at fixed $N$, and then
using the diagonal bound proves the first estimate in
\eqref{eq:boundary-average-variance}. Minkowski's inequality gives
the second from \eqref{eq:depth-moment-bound}, uniformly in $N$.
For $p\le2$ use the $L^2$ bound; for $p>2$ interpolate $L^2$ with
an $L^q$ bound for any $q>p$. A diagonal choice of finite depths as
$N$ increases gives the claimed reference sequence. In particular,
we have not asserted that one comparable-depth average, or a fixed
finite collection of them, becomes deterministic.
\end{proof}

\subsection{Absolute moments and the absence of diagonal terms}

Define the distributional moment $T_j$ by
\begin{equation}
 \langle T_j,f_1\otimes\cdots\otimes f_j\rangle
       =\EE\prod_{i=1}^j H(f_i),\qquad T_0=1.
 \label{eq:moment-distributions}
\end{equation}
These are well-defined distributions on $D^j$. For example, the
local Sobolev moment bounds in the proof of \cref{prop:tightness}
permit taking the expectation of the tensor power as a distribution;
equivalently, the continuous multilinear form in
\eqref{eq:moment-distributions} has a distributional kernel. Let
\[
 \mathcal C_j=\{(x_1,\ldots,x_j)\in D^j:x_i\ne x_l\text{ for }i\ne l\}
\]
be the configuration space without collisions.

\begin{proposition}[Harmonicity of absolute-height moments]
\label{prop:moment-harmonicity}
On $\mathcal C_j$, $T_j$ is a smooth function and is harmonic in each
variable separately.
\end{proposition}

\begin{proof}
The identity already transferred to the true law says
\begin{equation}
                     \EE[H(\Delta\phi)F]=0
 \label{eq:true-local-increment-identity}
\end{equation}
for bounded products of localized increment functions satisfying the
separation conditions of \cref{prop:bracket-identity}. In particular,
each factor may be a bounded continuous function of $H(q)$, where
$q$ has integral zero and is supported in a small disk. Moment bounds
and truncation extend the identity to products of finitely many such
$H(q)$'s whenever the same geometric conditions hold. The disks may
overlap one another.

Fix distinct points $x_0,x_2,\ldots,x_j$. Take a small closed disk
about $x_0$ with a larger concentric neighborhood compactly contained
in $D$ and disjoint from the spectator points. Choose spectator
neighborhoods and tests $f_i$ therein, and write $c_i=\int f_i$.
All these neighborhoods can be made as small as needed below. We
will approximate each absolute observation $H(f_i)$ by constants
and local zero-total observations supported a fixed positive distance
from the Laplace insertion.

Choose a unit smooth average $\rho_i$ near $x_i$, and a clean boundary
subarc accessible without entering a fixed larger disk about $x_0$.
Such access is always possible. A path in $D$ that meets the excluded
small disk can be detoured through a concentric annulus compactly
contained in $D$; thus removing this disk does not block a route to
the boundary subarc. For finitely many reference endpoints in
\cref{lem:boundary-references}, join $\rho_i$ to each endpoint by a
finite chain of unit smooth averages
\[
 \rho_{i,n,0}=\rho_i,\ \rho_{i,n,1},\ldots,
 \rho_{i,n,L_{i,n}}=\rho_{x_n(t),\kappa d(x_n(t))}.
\]
Consecutive supports lie in one disk of arbitrarily small prescribed
maximum diameter. Construct this by a sufficiently fine subdivision
of a path in its open tube, reducing the radii as necessary near the
boundary. All chain disks avoid the fixed neighborhood of $x_0$;
different chains need not avoid one another. There is the exact
identity
\begin{align}
 H(f_i)={}&c_i k+H(f_i-c_i\rho_i)
   +\frac{c_i}{N}\sum_{n=1}^N\sum_{l=0}^{L_{i,n}-1}
                 H(\rho_{i,n,l}-\rho_{i,n,l+1})
   +c_i(A_{N,t}-k).
 \label{eq:transport-identity}
\end{align}
Every nonconstant term before the last has total integral zero and
is supported in one small disk. The final term tends to zero in
every required $L^p$, with depth first and then $N$ increasing.

We check the separation geometry uniformly in the number of chain
links. Translate coordinates so that $x_0=0$, and suppose all chain
disks are outside $B(0,R)$. Choose $\ell\ll R$, and support $\phi$
in $B(0,\ell/10)$. Available lower and upper levels lie in
$(-2\ell,-\ell)$ and $(\ell,2\ell)$, respectively. Cap the diameter
of every chain disk, and of each initial spectator neighborhood,
by $\ell/(100j)$. Expand a product of the approximations on the
right side of \eqref{eq:transport-identity}, omitting the last error
term. Each expanded term contains at most $j-1$ disks: it selects
at most one local observation from each original spectator factor.
Thus the total length of the forbidden vertical projections is
less than $\ell/100$ in either available level interval, regardless
of the number or length of the routes. Choose regular levels from
the dense union of grids outside these projections and with positive
clearance.

Any selected disk in the resulting middle band is outside $B(0,R)$
and has vertical coordinate of size at most $2\ell$. It consequently
lies well to the left or right of the insertion, also in projection
along $u=(\sqrt3/2,1/2)$, because $\ell\ll R$. The boundary pieces
in the middle band have the same property by the interior clearance
of $x_0$. These are exactly the geometric conditions of
\cref{prop:bracket-identity}. Each finite expanded term therefore
satisfies \eqref{eq:true-local-increment-identity}; the levels can
depend on that term. There is no requirement that a single pair
avoid every disk in all the routes at once.

For each finite approximation the expanded sum is finite, so its
expectation against $H(\Delta\phi)$ is zero. The approximated
spectator variable differs from $H(f_i)$ by $c_i(A_{N,t}-k)$; in
particular its higher moments stay bounded independently of the
number of links. H\"older's inequality now passes to the $L^p$
reference limit in the unexpanded product. We obtain
\[
             \EE\bigl[H(\Delta\phi)\prod_{i=2}^jH(f_i)\bigr]=0.
\]
This proves distributional harmonicity in the distinguished variable
on a neighborhood of every collision-free configuration. Product
tests suffice by density in the smooth tests on a product
neighborhood. Repeating with each distinguished variable gives
$\Delta_{x_i}T_j=0$ on $\mathcal C_j$. Their sum is the ordinary
elliptic Laplacian on $\R^{2j}$, so distributional harmonic
regularity gives a smooth representative there. The separate
identities then hold classically.
\end{proof}

\begin{lemma}[No distributions supported on collision diagonals]
\label{lem:no-contact}
The smooth function $T_j$ on $\mathcal C_j$ is locally integrable on
$D^j$ and represents its full moment distribution. For $j\ge2$,
$K\Subset D$, and distinct $x_1,\ldots,x_j\in K$, with
$d_* =\min_{i<l}|x_i-x_l|$, it satisfies
\begin{equation}
 |T_j(x_1,\ldots,x_j)|
       \le C_{K,j}\bigl(1+\log_+(1/d_*)\bigr)^{j/2}.
 \label{eq:offdiagonal-log-bound}
\end{equation}
With distinct spectators $y_2,\ldots,y_j$ fixed in the interior,
$T_j(x,y_2,\ldots,y_j)$ is bounded as $x$ approaches any part of
$\partial D$, including the marked points. At a nonmarked point of
an arc of value $k$ its trace is
\begin{equation}
 \lim_{x\to p}T_j(x,y_2,\ldots,y_j)
                  =k\,T_{j-1}(y_2,\ldots,y_j).
 \label{eq:moment-boundary-trace}
\end{equation}
\end{lemma}

\begin{proof}
Radially mollify every variable at the common radius $r$ and denote
the result by $T_j^{(r)}$. For centers in $K$ and sufficiently small
$r$, the smearing bounds and H\"older give
\begin{equation}
 |T_j^{(r)}(x_1,\ldots,x_j)|
  =\left|\EE\prod_{i=1}^j H(\rho_{x_i,r})\right|
       \le C_{K,j}(1+|\log r|)^{j/2},
 \label{eq:common-mollifier-bound}
\end{equation}
even if some centers coincide. When $d_*>2r$, the averaging balls
are disjoint and separate harmonicity gives
$T_j^{(r)}=T_j$ exactly by successive applications of the mean-value
property. Taking $r$ to be a fixed small fraction of $d_*$, with a
fixed upper cutoff depending on $K$, proves
\eqref{eq:offdiagonal-log-bound}. Its right side is integrable on
$K^j$: it is bounded by a constant times one plus the sum over
pairs of powers of $|\log|x_i-x_l||$, and such powers are
integrable in two dimensions.

The collision neighborhood $\{d_*\le2r\}\cap K^j$ has volume
$O_{K,j}(r^2)$. By \eqref{eq:common-mollifier-bound}, the integral
of $|T_j^{(r)}|$ over that neighborhood tends to zero. The integral
there of the off-diagonal function tends to zero by its local
integrability. Off that neighborhood the two functions agree.
Thus $T_j^{(r)}$ converges in $L^1$ on compact sets to the extended
off-diagonal function. It also converges distributionally to the
original $T_j$, by the approximate-identity property. The two are
therefore the same distribution; no additional term supported on
a diagonal remains.

For the boundary claims, choose fixed, disjoint, small radial unit
averages $\rho_i$ at the spectator points. For $x$ sufficiently
near the boundary, its comparable-depth averaging ball is disjoint
from all of them. Separate mean value gives
\[
 T_j(x,y_2,\ldots,y_j)
   =\EE\left[H(\rho_{x,\kappa d(x)})
                         \prod_{i=2}^jH(\rho_i)\right].
\]
H\"older and \eqref{eq:depth-moment-bound} bound this uniformly in
$x$, including approaches to the marks. At a clean point,
\eqref{eq:mixed-boundary-trace} gives its limit as
$k\EE\prod_{i=2}^jH(\rho_i)$, which equals the right side of
\eqref{eq:moment-boundary-trace} by mean value once more.
\end{proof}
